\chapter{Fazit}
In diesem Versuch wurde die Compton-Streuung von $\gamma$-Strahlung an Aluminium experimentell untersucht. 
Hierzu wurde eine $^{137}$Cs-Quelle verwendet, deren Photonen auf Aluminiumtargets unterschiedlicher Dicke trafen. 
Zur Energiekalibrierung des Spektrometers diente zusätzlich eine $^{152}$Eu-Quelle mit bekannten $\gamma$-Übergängen.

Die Detektion der gestreuten beziehungsweise abgeschwächten Strahlung erfolgte mithilfe eines NaI(Tl)-Szintillationsdetektors mit angeschlossenem Photomultiplier. 
Die Signale wurden elektronisch verstärkt und anschließend mit einem Vielkanalanalysator in Energiespektren umgewandelt.

Im ersten Teil des Versuchs wurde die Abschwächung der $\gamma$-Strahlung in Aluminium untersucht. 
Die gemessene Transmission zeigte insgesamt den erwarteten exponentiellen Verlauf gemäß dem Lambert-Beer-Gesetz. 
Lediglich ein Messpunkt wich aufgrund geometrischer Eigenschaften des verwendeten Targets deutlich vom theoretischen Verhalten ab.

Aus den Transmissionsmessungen wurde ein totaler Wirkungsquerschnitt von
\begin{align}
    \sigma_\mathrm{tot}
    =
    (188,6 \pm 5,7)\,\mathrm{mb}
\end{align}
bestimmt.

Dieser Wert liegt in derselben Größenordnung wie der theoretische Wirkungsquerschnitt nach der Klein-Nishina-Theorie sowie der Literaturwert:
\begin{align}
    \sigma_{\mathrm{tot,KN}}
    &\approx
    256\,\mathrm{mb},
    \\
    \sigma_{\mathrm{tot,lit}}
    &\approx
    257\,\mathrm{mb}.
\end{align}
Die verbleibende Abweichung kann insbesondere auf systematische Unsicherheiten zurückgeführt werden, beispielsweise auf Untergrundbeiträge, geometrische Effekte des Aufbaus oder Unsicherheiten bei der Effizienzkorrektur des Detektors.

Anschließend wurde das Spektrometer energiekalibriert. 
Dazu wurde die Cs-Quelle abgeschirmt und eine Eu-Quelle in den Strahlengang eingefügt. 
Das aufgezeichnete Spektrum zeigte deutlich mehrere Peaks, die den bekannten Übergangsenergien des Europium-Zerfalls entsprachen.
Eine lineare Regression der Kanalnummern gegen die in der Literatur angegebenen Energien diente als Kalibrierfunktion und bestätigte den linearen Betrieb des Vielkanalanalysators. 
Durch Vergleich der relativen Intensitäten der Spektrallinien mit bekannten Übergangswahrscheinlichkeiten konnten zudem einige Punkte der relativen Nachweiseffizienz des NaI(Tl)-Detektors bestimmt werden. 
Aufgrund der begrenzten Anzahl verfügbarer Datenpunkte lieferte die mit einem linearen kontinuierlichen Untergrund angepasste Gauß-Funktion jedoch nur eine eingeschränkte Zuverlässigkeit für die Effizienz über den gesamten Energiebereich.

In der letzten Messreihe wurde das Compton-Streuspektrum einer $^{137}\text{Cs}$-Quelle mit einem Target fester Dicke von 1 mm und verschiedenen Streuwinkeln aufgenommen. 
Für jede Messung wurde ein kontinuierliches Hintergrundbild ohne Target erstellt und vom Spektrum subtrahiert. 
Die Analyse der erhaltenen Spektren ermöglichte eine detaillierte Überprüfung der Compton-Gleichung : Den im Spektrum identifizierten Maxima wurden theoretische Energien zugeordnet und grafisch als Funktion des Streuwinkels dargestellt. 
Die Anpassung ergab eine einfallende Photonenenergie von:$E_{\gamma} = (637.86 \pm7.35)$ keV, welcher um $4\%$ von dem in der Literatur angegebene Wert von $661,7keV$ und stimmt innerhalb einer $3\sigma$-Abweichung mit diesem überein. 

Zusätzlich wurde die Winkelabhängigkeit der gestreuten Photonenenergie untersucht. 
Die gemessenen Energien folgen insgesamt der erwarteten Compton-Beziehung und bestätigen die theoretisch vorhergesagte Abnahme der Photonenergie mit wachsendem Streuwinkel.

Darüber hinaus wurde die Winkelabhängigkeit der Streuintensität mit der Klein-Nishina-Theorie verglichen. 
Nach Durchführung der Effizienzkorrektur zeigte sich eine qualitative Übereinstimmung mit dem theoretisch erwarteten Verlauf. 
Insbesondere die Abnahme der Intensität bei großen Streuwinkeln konnte experimentell nachvollzogen werden.

Insgesamt bestätigen die Messergebnisse die wesentlichen Vorhersagen der Compton-Theorie sowie die grundlegende Winkel- und Energieabhängigkeit der Compton-Streuung.

Zusammenfassend lässt sich sagen, dass der gesamte Compton-Streuquerschnitt der $^{137}\text{Cs}$-Quelle quantitativ und die Winkelabhängigkeit der Energie und Intensität der gestreuten Photonen qualitativ bestimmt werden konnten.
